\chapter{2004年浙江卷（理）}
\section{单选题}
\begin{problem}
若 \(U=\{1,2,3,4\}\)，\(M=\{1,2\}\)，\(N=\{2,3\}\)，则 \(\complement_U(M\cup N)=\)（\quad）

\choices
    {\(\{1,2,3\}\)}
    {\(\{2\}\)}
    {\(\{1,3,4\}\)}
    {\(\{4\}\)}
\end{problem}

\begin{answer}
D
\end{answer}

\begin{solution}
先求并集：\(M\cup N=\{1,2,3\}\). 补集中的元素必须属于全集 \(U\)，但不属于这个并集，因此
\(\complement_U(M\cup N)=\{4\}\)，故选 D.
\end{solution}

\begin{problem}
点 \(P\) 从 \((1,0)\) 出发，沿单位圆 \(x^2+y^2=1\) 逆时针方向运动 \(\frac{2\pi}{3}\) 弧长到达 \(Q\) 点，则 \(Q\) 的坐标为（\quad）

\choices
    {\(\left(-\frac{1}{2}, \frac{\sqrt{3}}{2}\right)\)}
    {\(\left(-\frac{\sqrt{3}}{2}, -\frac{1}{2}\right)\)}
    {\(\left(-\frac{1}{2}, -\frac{\sqrt{3}}{2}\right)\)}
    {\(\left(-\frac{\sqrt{3}}{2}, \frac{1}{2}\right)\)}
\end{problem}

\begin{answer}
A
\end{answer}

\begin{solution}
单位圆的半径为 \(r=1\). 由弧长公式 \(s=r\theta\)，点 \(P\) 对应的圆心角为
\(\theta=\frac{s}{r}=\frac{2\pi/3}{1}=\frac{2\pi}{3}\). 从 \((1,0)\) 沿逆时针方向转过 \(\frac{2\pi}{3}\) 后，
\[
Q=\left(\cos\frac{2\pi}{3},\sin\frac{2\pi}{3}\right)=\left(-\frac{1}{2},\frac{\sqrt{3}}{2}\right)
\]
故选 A.
\end{solution}

\begin{problem}
已知等差数列 \(\{a_n\}\) 的公差为 2，若 \(a_1\)，\(a_3\)，\(a_4\) 成等比数列，则 \(a_2=\)（\quad）

\choices
    {\(-4\)}
    {\(-6\)}
    {\(-8\)}
    {\(-10\)}
\end{problem}

\begin{answer}
B
\end{answer}

\begin{solution}
设首项为 \(a_1\). 由公差为 \(2\)，
\(a_3=a_1+2\times2=a_1+4\)，\(a_4=a_1+3\times2=a_1+6\).
三项 \(a_1\)，\(a_3\)，\(a_4\) 成等比数列，所以中项平方等于两端项之积，即
\[
(a_1+4)^2=a_1(a_1+6)
\]
展开并约去 \(a_1^2\)，得 \(8a_1+16=6a_1\)，从而 \(a_1=-8\). 因此
\(a_2=a_1+2=-6\)，故选 B.
\end{solution}

\begin{problem}
曲线 \(y^2=4x\) 关于直线 \(x=2\) 对称的曲线方程是（\quad）

\choices
    {\(y^2=8-4x\)}
    {\(y^2=4x-8\)}
    {\(y^2=16-4x\)}
    {\(y^2=4x-16\)}
\end{problem}

\begin{answer}
C
\end{answer}

\begin{solution}
设原曲线上任一点为 \((x,y)\). 关于直线 \(x=2\) 对称后，纵坐标不变，横坐标变为
\(X=2\times2-x=4-x\)，所以 \(x=4-X\). 代入原方程得
\[
y^2=4x=4(4-X)=16-4X
\]
把对称后坐标的字母 \(X\) 改写为通常的 \(x\)，所得曲线方程为 \(y^2=16-4x\)，故选 C.
\end{solution}

\begin{problem}
设 \(z=x-y\)，式中变量 \(x\) 和 \(y\) 满足条件 \(\begin{cases}
x+y-3\geqslant0,\\
x-2y\geqslant0,
\end{cases}\) 则 \(z\) 的最小值为（\quad）

\choices
    {1}
    {\(-1\)}
    {3}
    {\(-3\)}
\end{problem}

\begin{answer}
A
\end{answer}

\begin{solution}
令 \(u=x+y-3\)，\(v=x-2y\)，则 \(u\geqslant0\)，\(v\geqslant0\)，且
\(x+y=3+u\)，\(x-2y=v\).
两式相减消去 \(x\)，得 \(3y=3+u-v\)，所以
\(y=1+\frac{u-v}{3}\)，进而 \(x=2+\frac{2u+v}{3}\). 因此
\[
z=x-y=1+\frac{u+2v}{3}\geqslant1
\]
当 \(u=v=0\)，即 \((x,y)=(2,1)\) 时满足原条件且取等号，所以 \(z\) 的最小值为 \(1\)，故选 A.
\end{solution}

\begin{problem}
已知复数 \(z_1=3+4\i\)，\(z_2=t+\i\)，且 \(z_1\cdot z_2\) 为实数，则实数 \(t=\)（\quad）

\choices
    {\(\frac{3}{4}\)}
    {\(\frac{4}{3}\)}
    {\(-\frac{4}{3}\)}
    {\(-\frac{3}{4}\)}
\end{problem}

\begin{answer}
D
\end{answer}

\begin{solution}
利用 \(\i^2=-1\)，
\[
z_1z_2=(3+4\i)(t+\i)=3t+3\i+4t\i+4\i^2=(3t-4)+(3+4t)\i
\]
复数为实数的充要条件是虚部为零，因此 \(3+4t=0\)，解得
\(t=-\frac34\)，故选 D.
\end{solution}

\begin{problem}
若 \(\left(\sqrt{x}+\frac{2}{\sqrt[3]{x}}\right)^n\) 展开式中存在常数项，则 \(n\) 的值可以是（\quad）

\choices
    {8}
    {9}
    {10}
    {12}
\end{problem}

\begin{answer}
C
\end{answer}

\begin{solution}
在展开式的一般项中，若从第二项取 \(k\) 次，则从第一项取 \(n-k\) 次，其中 \(0\leqslant k\leqslant n\). 该项中 \(x\) 的指数为
\[
\frac{n-k}{2}-\frac{k}{3}=\frac{3n-5k}{6}
\]
常数项要求指数为零，所以 \(3n=5k\). 因为 \(3\) 与 \(5\) 互质，故 \(5\mid n\). 给出的正整数中只有 \(n=10\) 是 \(5\) 的倍数；此时 \(k=\frac{3n}{5}=6\)，满足 \(0\leqslant k\leqslant n\)，故选 C.
\end{solution}

\begin{problem}
在 \(\triangle ABC\) 中，“\(A>30^{\circ}\)”是“\(\sin A>\frac{1}{2}\)”的（\quad）

\choices
    {充分而不必要条件}
    {必要而不充分条件}
    {充分必要条件}
    {既不充分也必要条件}
\end{problem}

\begin{answer}
B
\end{answer}

\begin{solution}
三角形内角满足 \(0<A<\pi\). 在此范围内，
\[
\sin A>\frac12\iff \frac{\pi}{6}<A<\frac{5\pi}{6}
\]
所以 \(\sin A>\frac12\) 一定能推出 \(A>\frac{\pi}{6}=30^{\circ}\)，即“\(A>30^{\circ}\)”是“\(\sin A>\frac12\)”的必要条件. 反过来不成立，例如取 \(A=\frac{5\pi}{6}=150^{\circ}\)，并令另外两个角均为 \(15^{\circ}\)，仍构成三角形，但 \(\sin A=\frac12\) 不满足严格不等式. 因此“\(A>30^{\circ}\)”是“\(\sin A>\frac12\)”的必要而不充分条件，故选 B.
\end{solution}

\begin{problem}
若椭圆 \(\frac{x^2}{a^2}+\frac{y^2}{b^2}=1\)（\(a>b>0\)）的左、右焦点分别为 \(F_1\)，\(F_2\)，线段 \(F_1F_2\) 被抛物线 \(y^2=2bx\) 的焦点分成 \(5:3\) 两段，则此椭圆的离心率为（\quad）

\choices
    {\(\frac{16}{17}\)}
    {\(\frac{4\sqrt{17}}{17}\)}
    {\(\frac{4}{5}\)}
    {\(\frac{2\sqrt{5}}{5}\)}
\end{problem}

\begin{answer}
D
\end{answer}

\begin{solution}
设椭圆的半焦距为 \(c\)，则焦点为 \(F_1(-c,0)\)，\(F_2(c,0)\). 将抛物线写成标准形式
\(y^2=4px\)，由 \(4p=2b\) 得 \(p=\frac b2\)，所以其焦点为 \(S\left(\frac b2,0\right)\).
由于 \(S\) 在线段 \(F_1F_2\) 上且将其分成 \(5:3\) 两段，
\[
\frac{SF_1}{SF_2}=\frac{c+b/2}{c-b/2}=\frac53
\]
于是 \(3c+\frac32b=5c-\frac52b\)，即 \(c=2b\). 椭圆参数满足
\(a^2=b^2+c^2\)，故 \(a^2=5b^2\). 因而离心率为
\[
e=\frac ca=\frac{2b}{\sqrt5b}=\frac2{\sqrt5}=\frac{2\sqrt5}{5}
\]
故选 D.
\end{solution}

\begin{problem}
如图，在正三棱柱 \(ABC-A_{1}B_{1}C_{1}\) 中已知 \(AB=1\)，\(D\) 在棱 \(BB_{1}\) 上，且 \(BD=1\)，若 \(AD\) 与平面 \(AA_{1}C_{1}C\) 所成的角为 \(\alpha\)，则 \(\alpha=\)（\quad）

\bitmapfigure[width=0.32\linewidth]{img/2004/zhejiang_science/q10_fig1.png}

\choices
    {\(\frac{\pi}{3}\)}
    {\(\frac{\pi}{4}\)}
    {\(\arcsin \frac{\sqrt{10}}{4}\)}
    {\(\arcsin \frac{\sqrt{6}}{4}\)}
\end{problem}

\begin{answer}
D
\end{answer}

\begin{solution}
建立直角坐标系，使底面在 \(xOy\) 平面内，取
\(A=(0,0,0)\)，\(B=(1,0,0)\)，\(C=\left(\frac12,\frac{\sqrt3}{2},0\right)\).
因为侧棱垂直于底面且 \(BD=1\)，可取 \(D=(1,0,1)\). 平面 \(AA_1C_1C\) 由方向向量
\(\overrightarrow{AC}=\left(\frac12,\frac{\sqrt3}{2},0\right)\) 和竖直方向共同确定，因而可取其单位法向量
\(\bs n=\left(\frac{\sqrt3}{2},-\frac12,0\right)\). 又
\(\overrightarrow{AD}=(1,0,1)\)，所以
\(|\overrightarrow{AD}\cdot\bs n|=\frac{\sqrt3}{2}\)，\(|AD|=\sqrt{1^2+1^2}=\sqrt2\).
直线与平面所成角为 \(\alpha\) 时，\(\sin\alpha\) 等于方向向量在平面法向上的投影长度与方向向量长度之比，故
\[
\sin\alpha=\frac{\sqrt3/2}{\sqrt2}=\frac{\sqrt6}{4}
\]
因此 \(\alpha=\arcsin\frac{\sqrt6}{4}\)，故选 D.
\end{solution}

\begin{problem}
设 \(f'(x)\) 是函数 \(f(x)\) 的导函数，\(y=f'(x)\) 的图象如图所示，则 \(y=f(x)\) 的图象最有可能的是（\quad）

\bitmapfigure[width=0.42\linewidth]{img/2004/zhejiang_science/q11_fig1.png}

\choices
    {\choicebitmap[width=0.15\paperwidth]{img/2004/zhejiang_science/q11_optA.png}}
    {\choicebitmap[width=0.15\paperwidth]{img/2004/zhejiang_science/q11_optB.png}}
    {\choicebitmap[width=0.15\paperwidth]{img/2004/zhejiang_science/q11_optC.png}}
    {\choicebitmap[width=0.15\paperwidth]{img/2004/zhejiang_science/q11_optD.png}}
\end{problem}

\begin{answer}
C
\end{answer}

\begin{solution}
由图可读出 \(f'(0)=f'(2)=0\)，并且
\(f'(x)>0\)，\((x<0\text{ 或 }x>2)\)，\(f'(x)<0\)，\((0<x<2)\).
所以 \(f(x)\) 在 \((-\infty,0)\) 上递增，在 \((0,2)\) 上递减，在 \((2,+\infty)\) 上递增. 于是 \(x=0\) 处由增变减，是极大值点；\(x=2\) 处由减变增，是极小值点. 四个选项中只有 C 的图象同时具有这两个单调性特征，故选 C.
\end{solution}

\begin{problem}
若 \(f(x)\) 和 \(g(x)\) 都是定义在实数集 \(\R\) 上的函数，且方程 \(x-f[g(x)]=0\) 有实数解，则 \(g[f(x)]\) 不可能是（\quad）

\choices
    {\(x^2+x-\frac{1}{5}\)}
    {\(x^2+x+\frac{1}{5}\)}
    {\(x^2-\frac{1}{5}\)}
    {\(x^2+\frac{1}{5}\)}
\end{problem}

\begin{answer}
B
\end{answer}

\begin{solution}
设方程 \(x-f(g(x))=0\) 的实数解为 \(x_0\)，则 \(x_0=f(g(x_0))\). 令 \(y_0=g(x_0)\)，便有 \(x_0=f(y_0)\)，从而
\[
g(f(y_0))=g(x_0)=y_0
\]
因此函数 \(g\circ f\) 必须有实数不动点，即方程 \(g(f(x))=x\) 必须有实数解.

逐项判断：
\[
x^2+x-\frac15=x\iff x^2-\frac15=0
\]
有实数解；
\[
x^2+x+\frac15=x\iff x^2+\frac15=0
\]
无实数解；
\[
x^2-\frac15=x
\]
的判别式为 \(1+\frac45>0\)，有实数解；
\[
x^2+\frac15=x
\]
的判别式为 \(1-\frac45>0\)，也有实数解. 因此不可能的是 B.
\end{solution}

\section{填空题}
\begin{problem}
已知 \(f(x)=\begin{cases}
1,&x\geqslant0,\\
-1,&x<0,
\end{cases}\) 则不等式 \(x+(x+2)\cdot f(x+2)\leqslant5\) 的解集是\fillinblank{}.
\end{problem}

\begin{answer}
\((-\infty,\frac{3}{2}]\)
\end{answer}

\begin{solution}
分界点由 \(x+2=0\) 给出，即 \(x=-2\). 当 \(x\geqslant-2\) 时，
\(f(x+2)=1\)，原不等式化为
\[
x+(x+2)\leqslant5\iff 2x\leqslant3\iff x\leqslant\frac32
\]
结合本段条件，得到 \(\left[-2,\frac32\right]\).

当 \(x<-2\) 时，\(f(x+2)=-1\)，左端为
\[
x-(x+2)=-2\leqslant5
\]
恒成立，所以得到 \((-\infty,-2)\). 合并两段解集，且 \(x=-2\) 已包含在第一段中，故解集为
\(\left(-\infty,\frac32\right]\).
\end{solution}

\begin{problem}
已知平面上三点 \(A\)，\(B\)，\(C\) 满足 \(\left|\overrightarrow{AB}\right|=3\)，\(\left|\overrightarrow{BC}\right|=4\)，\(\left|\overrightarrow{CA}\right|=5\)，则 \(\overrightarrow{AB}\cdot\overrightarrow{BC}+\overrightarrow{BC}\cdot\overrightarrow{CA}+\overrightarrow{CA}\cdot\overrightarrow{AB}\) 的值等于\fillinblank{}.
\end{problem}

\begin{answer}
\(-25\)
\end{answer}

\begin{solution}
记题目所求的量为 \(T\). 首尾相接的向量满足
\(\overrightarrow{AB}+\overrightarrow{BC}+\overrightarrow{CA}=\bs{0}\). 两边取模平方，得
\[
0=|\overrightarrow{AB}|^2+|\overrightarrow{BC}|^2+|\overrightarrow{CA}|^2
  +2\left(\overrightarrow{AB}\cdot\overrightarrow{BC}
  +\overrightarrow{BC}\cdot\overrightarrow{CA}
  +\overrightarrow{CA}\cdot\overrightarrow{AB}\right)
\]
代入三条向量的模，并用 \(T\) 表示括号中的三项和，得
\[
0=3^2+4^2+5^2+2T=50+2T
\]
所以 \(T=-25\).
\end{solution}

\begin{problem}
设坐标平面内有一个质点从原点出发，沿 \(x\) 轴跳动. 每次向正方向或负方向跳 1 个单位，经过 5 次跳动质点落在点 \((3,0)\)（允许重复过此点）处，则质点不同的运动方法共有\fillinblank{} 种.（用数字作答）
\end{problem}

\begin{answer}
5
\end{answer}

\begin{solution}
设向正方向跳的次数为 \(p\)，向负方向跳的次数为 \(q\). 总跳动次数和最终位移分别给出
\(p+q=5\)，\(p-q=3\).
两式相加、相减得 \(p=4\)，\(q=1\). 因而五次跳动中恰有一次为负方向跳动；确定这一次发生在第几次即可唯一确定整个运动方法，所以方法数为
\[
\mathrm C_5^1=5
\]
\end{solution}

\begin{problem}
已知平面 \(\alpha\) 和平面 \(\beta\) 交于直线 \(l\)，\(P\) 是空间一点，\(PA\perp \alpha\)，垂足为 \(A\)，\(PB\perp \beta\)，垂足为 \(B\)，且 \(PA=1\)，\(PB=2\)，若点 \(A\) 在 \(\beta\) 内的射影与点 \(B\) 在 \(\alpha\) 内的射影重合，则点 \(P\) 到 \(l\) 的距离为\fillinblank{}.
\end{problem}

\begin{answer}
\(\sqrt{5}\)
\end{answer}

\begin{solution}
设两个射影的重合点为 \(C\). 由 \(AC\perp\beta\)，且 \(A\)，\(C\in\alpha\)，若 \(A\ne C\)，则平面 \(\alpha\) 含有一条垂直于平面 \(\beta\) 的直线，所以 \(\alpha\perp\beta\). 若 \(A=C\)，则 \(C\in l\). 又 \(PA=1\ne2=PB\)，故 \(B\ne C\)；此时 \(B\)，\(C\in\beta\) 且 \(BC\perp\alpha\)，同样得到 \(\alpha\perp\beta\).

建立直角坐标系，使交线 \(l\) 为 \(z\) 轴，\(\alpha\) 为 \(xOz\) 平面，\(\beta\) 为 \(yOz\) 平面. 设 \(P=(x,y,z)\)，则 \(P\) 到 \(\alpha\) 的垂距为 \(|y|\)，到 \(\beta\) 的垂距为 \(|x|\)，所以
\(|y|=PA=1\)，\(|x|=PB=2\).
点 \(P\) 到 \(z\) 轴，即到 \(l\) 的距离为
\[
\sqrt{x^2+y^2}=\sqrt{2^2+1^2}=\sqrt5
\]
\end{solution}

\section{解答题}
\begin{problem}
在 \(\triangle ABC\) 中，角 \(A\)，\(B\)，\(C\) 所对的边分别为 \(a\)，\(b\)，\(c\)，且 \(\cos A=\frac{1}{3}\).
\begin{enumerate}
\item 求 \(\sin^{2}\frac{B+C}{2}+\cos 2A\) 的值；
\item 若 \(a=\sqrt{3}\)，求 \(bc\) 的最大值.
\end{enumerate}
\end{problem}

\begin{solution}
\begin{enumerate}
\item 因为 \(A+B+C=\pi\)，所以 \(B+C=\pi-A\)，从而
\[
\sin^2\frac{B+C}{2}
=\sin^2\left(\frac{\pi}{2}-\frac A2\right)
=\cos^2\frac A2
=\frac{1+\cos A}{2}
=\frac23
\]
另外，
\[
\cos 2A=2\cos^2A-1=2\left(\frac13\right)^2-1=-\frac79
\]
因此所求的值为
\(\frac23-\frac79=-\frac19\).
\item 由余弦定理及 \(\cos A=\frac13\)，且 \(a=\sqrt3\)，得
\[
3=a^2=b^2+c^2-2bc\cos A=b^2+c^2-\frac23bc
=(b-c)^2+\frac43bc
\]
所以
\[
bc=\frac34\left(3-(b-c)^2\right)\leqslant\frac94
\]
当且仅当 \(b=c\) 时等号成立. 代入上式，此时
\(3=\frac43b^2\)，得 \(b=c=\frac32\)，且三角形条件成立. 因此 \(bc\) 的最大值为 \(\frac94\).
\end{enumerate}
\end{solution}

\begin{problem}
盒子中有大小相同的球 10 个，其中标号为 1 的球 3 个，标号为 2 的球 4 个，标号为 5 的球 3 个，第一次从盒子中任取 1 个球，放回后第二次再任取 1 个球（假设取到每个球的可能性都相同）. 记第一次与第二次取到球的标号之和为 \(\varepsilon\).
\begin{enumerate}
\item 求随机变量 \(\varepsilon\) 的分布列；
\item 求随机变量 \(\varepsilon\) 的期望 \(E(\varepsilon)\).
\end{enumerate}
\end{problem}

\begin{solution}
\begin{enumerate}
\item 每次取到标号 \(1\)，\(2\)，\(5\) 的概率分别为
\(\frac3{10}\)，\(\frac4{10}\)，\(\frac3{10}\). 由于取球后放回，两次取球相互独立. 设两次取到的标号分别为 \(\varepsilon_1\)，\(\varepsilon_2\)，则
\(\varepsilon=\varepsilon_1+\varepsilon_2\) 的可能取值为 \(2\)，\(3\)，\(4\)，\(6\)，\(7\)，\(10\). 逐一按有序组合计算：
\(P(\varepsilon=2)=\frac3{10}\cdot\frac3{10}=\frac9{100}\)，\(P(\varepsilon=3)=2\cdot\frac3{10}\cdot\frac4{10}=\frac{24}{100}\)，\(P(\varepsilon=4)=\frac4{10}\cdot\frac4{10}=\frac{16}{100}\). 又 \(P(\varepsilon=6)=2\cdot\frac3{10}\cdot\frac3{10}=\frac{18}{100}\)，\(P(\varepsilon=7)=2\cdot\frac4{10}\cdot\frac3{10}=\frac{24}{100}\)，\(P(\varepsilon=10)=\frac3{10}\cdot\frac3{10}=\frac9{100}\).
这些概率之和为 \(1\)，故分布列为
\begin{center}
\begin{tabular}{c|cccccc}
\hline
\(\varepsilon\) & \(2\) & \(3\) & \(4\) & \(6\) & \(7\) & \(10\) \\
\hline
\(P\) & \(\frac9{100}\) & \(\frac{24}{100}\) & \(\frac{16}{100}\) & \(\frac{18}{100}\) & \(\frac{24}{100}\) & \(\frac9{100}\) \\
\hline
\end{tabular}
\end{center}
\item 对任意一次取球，标号的期望为
\[
E(\varepsilon_1)=E(\varepsilon_2)
=1\cdot\frac3{10}+2\cdot\frac4{10}+5\cdot\frac3{10}
=\frac{13}{5}
\]
由期望的线性性质，
\[
E(\varepsilon)=E(\varepsilon_1+\varepsilon_2)
=E(\varepsilon_1)+E(\varepsilon_2)
=\frac{26}{5}
\]
\end{enumerate}
\end{solution}

\begin{problem}
如图，已知正方形 \(ABCD\) 和矩形 \(ACEF\) 所在的平面互相垂直，\(AB=\sqrt{2}\)，\(AF=1\)，\(M\) 是线段 \(EF\) 的中点.
\begin{enumerate}
\item 求证：\(AM \parallel\) 平面 \(BDE\)；
\item 求二面角 \(A-DF-B\) 的大小.
\end{enumerate}

\FigureLayoutDeclare{img/2004/zhejiang_science/q19_fig1.png}{7.0cm}{95bp 77bp 23bp 11bp}
\bitmapfigure[width=0.32\linewidth]{img/2004/zhejiang_science/q19_fig1.png}
\end{problem}

\begin{solution}
\begin{enumerate}
\item 取直角坐标系，使 \(A=(-1,0,0)\)，\(C=(1,0,0)\)，正方形 \(ABCD\) 在 \(xOy\) 平面内，矩形 \(ACEF\) 在 \(xOz\) 平面内. 因为 \(AC=2\)，\(AB=\sqrt2\)，可取
\(B=(0,1,0)\)，\(D=(0,-1,0)\).
又 \(AF=1\)，从而
\(E=(1,0,1)\)，\(F=(-1,0,1)\)，\(M=(0,0,1)\).
\(\overrightarrow{AM}=(1,0,1)\)，\(\overrightarrow{BD}=(0,-2,0)\)，\(\overrightarrow{BE}=(1,-1,1)\).
平面 \(BDE\) 的一个法向量为
\[
\overrightarrow{BD}\times\overrightarrow{BE}=(-2,0,2)
\]
因此平面 \(BDE\) 的方程为 \(z=x\). 由于 \(\overrightarrow{AM}\cdot(-2,0,2)=0\)，直线 \(AM\) 的方向平行于平面 \(BDE\)；又 \(A=(-1,0,0)\) 不满足 \(z=x\)，所以 \(A\) 不在该平面上，从而 \(AM\parallel\) 平面 \(BDE\).
\item 平面 \(ADF\) 与平面 \(BDF\) 的法向量可分别取为
\(\bs n_1=\overrightarrow{DA}\times\overrightarrow{DF}=(1,1,0)\)，\(\bs n_2=\overrightarrow{DB}\times\overrightarrow{DF}=(2,0,2)\)
其中 \(\overrightarrow{DA}=(-1,1,0)\)，\(\overrightarrow{DB}=(0,2,0)\)，\(\overrightarrow{DF}=(-1,1,1)\). 所以二面角 \(A\mathord{-}DF\mathord{-}B\) 的余弦为
\[
\cos\theta=\frac{|\bs n_1\cdot\bs n_2|}{|\bs n_1||\bs n_2|}
=\frac{2}{\sqrt2\cdot2\sqrt2}=\frac12
\]
因此二面角 \(A\mathord{-}DF\mathord{-}B=\frac{\pi}{3}\).
\end{enumerate}
\end{solution}

\begin{problem}
设曲线 \(y=\e^{-x}\)（\(x\geqslant0\)）在点 \(M(t,\e^{-t})\) 处的切线 \(l\) 与 \(x\) 轴 \(y\) 轴所围成的三角形面积为 \(S(t)\).
\begin{enumerate}
\item 求切线 \(l\) 的方程；
\item 求 \(S(t)\) 的最大值.
\end{enumerate}
\end{problem}

\begin{solution}
\begin{enumerate}
\item 因为 \(y'=-\e^{-x}\)，所以在 \(M(t,\e^{-t})\) 处的切线斜率为 \(-\e^{-t}\). 由点斜式，切线 \(l\) 的方程为
\[
y-\e^{-t}=-\e^{-t}(x-t)
\]
整理得
\[
y=\e^{-t}(t+1-x)
\]
\item 令 \(y=0\)，得切线与 \(x\) 轴交于 \((t+1,0)\)；令 \(x=0\)，得其与 \(y\) 轴交于 \((0,(t+1)\e^{-t})\). 由于 \(t\geqslant0\)，两个截距均非负，三角形面积为
\[
S(t)=\frac12(t+1)^2\e^{-t}
\]
对 \(t\geqslant0\) 求导：
\[
S'(t)=\frac12\e^{-t}(t+1)(1-t)
\]
因为 \(\e^{-t}>0\)，\(t+1>0\)，所以 \(S'(t)>0\)（\(0\leqslant t<1\)），\(S'(1)=0\)，\(S'(t)<0\)（\(t>1\)）. 因而 \(S(t)\) 在 \([0,1]\) 上递增，在 \([1,+\infty)\) 上递减，最大值在 \(t=1\) 时取得，且
\[
S(1)=\frac12(1+1)^2\e^{-1}=\frac2{\e}
\]
\end{enumerate}
\end{solution}

\begin{problem}
已知双曲线的中心在原点，右顶点为 \(A(1,0)\)，点 \(P\)，\(Q\) 在双曲线的右支上，点 \(M(m,0)\) 到直线 \(AP\) 的距离为 1.
\begin{enumerate}
\item 若直线 \(AP\) 的斜率为 \(k\)，且 \(\left|k\right| \in \left[\frac{\sqrt{3}}{3}, \sqrt{3}\right]\)，求实数 \(m\) 的取值范围；
\item 当 \(m=\sqrt{2}+1\) 时，\(\triangle APQ\) 的内心恰好是点 \(M\)，求此双曲线的方程.
\end{enumerate}
\end{problem}

\begin{solution}
\begin{enumerate}
\item 直线 \(AP\) 的方程为 \(y=k(x-1)\)，即 \(kx-y-k=0\). 由点到直线的距离公式，
\[
\frac{|k(m-1)|}{\sqrt{1+k^2}}=1
\]
所以
\[
|m-1|=\frac{\sqrt{1+k^2}}{|k|}=\sqrt{1+\frac1{k^2}}
\]
由 \(\frac{\sqrt3}{3}\leqslant|k|\leqslant\sqrt3\)，得
\(\frac13\leqslant k^2\leqslant3\)，从而
\(\frac{2}{\sqrt3}\leqslant|m-1|\leqslant2\). 因此
\[
m\in\left[-1,1-\frac{2\sqrt3}{3}\right]\cup\left[1+\frac{2\sqrt3}{3},3\right]
\]
\item 由于双曲线的右顶点为 \(A(1,0)\)，其方程可设为
\(\frac{x^2}{1^2}-\frac{y^2}{b^2}=1\)，\(b>0\).
内心 \(M\) 在顶点 \(A\) 的内角平分线上，而 \(AM\) 是 \(x\) 轴，所以两条射线 \(AP\)，\(AQ\) 关于 \(x\) 轴对称. 对称的两条直线与该双曲线除 \(A\) 外的交点横坐标相同，故可设
\(P=(x,y)\)，\(Q=(x,-y)\)，\(y>0\).
于是 \(PQ\) 为直线 \(x=x\). 因为内心在三角形内部，\(x>m\)，且内心到边 \(PQ\) 的距离为 \(1\)，所以
\[
x-m=1
\]
代入 \(m=\sqrt2+1\)，得 \(x=2+\sqrt2\).

设直线 \(AP\) 的斜率为 \(k=\frac{y}{x-1}\). 由题设点到直线的距离为 \(1\)，以及 \(m-1=\sqrt2\)，有
\[
\frac{|k|\sqrt2}{\sqrt{1+k^2}}=1
\]
平方并整理得 \(2k^2=1+k^2\)，所以 \(k^2=1\). 由于 \(P\) 在上方且 \(x>1\)，取 \(k=1\)，从而
\[
y=k(x-1)=x-1=1+\sqrt2
\]
代入双曲线方程，得
\[
\frac1{b^2}=\frac{x^2-1}{y^2}
=\frac{(2+\sqrt2)^2-1}{(1+\sqrt2)^2}
=2\sqrt2-1
\]
故所求双曲线方程为
\[
x^2-(2\sqrt2-1)y^2=1
\]
\end{enumerate}
\end{solution}

\begin{problem}
如图，\(\triangle OBC\) 的三个顶点坐标分别为 \((0,0)\)，\((1,0)\)，\((0,2)\). 设 \(P_{1}\) 为线段 \(BC\) 的中点，\(P_{2}\) 为线段 \(CO\) 的中点，\(P_{3}\) 为线段 \(OP_{1}\) 的中点. 对于每一个正整数 \(n\)，\(P_{n+3}\) 为线段 \(P_{n}P_{n+1}\) 的中点，令 \(P_{n}\) 的坐标为 \((x_{n}, y_{n})\)，\(a_{n} = \frac{1}{2}y_{n} + y_{n+1} + y_{n+2}\).
\begin{enumerate}
\item 求 \(a_{1}\)，\(a_{2}\)，\(a_{3}\) 及 \(a_{n}\)；
\item 证明 \(y_{n+4} = 1 - \frac{y_n}{4}\)，\(n \in \N^*\)；
\item 若记 \(b_{n} = y_{4n+4} - y_{4n}\)，\(n \in \N^*\)，证明 \(\{b_n\}\) 是等比数列.
\end{enumerate}

\bitmapfigure[float=inline,width=0.32\linewidth]{img/2004/zhejiang_science/q22_fig1.png}
\end{problem}

\begin{solution}
\begin{enumerate}
\item 由中点坐标公式，
\(P_1=\left(\frac12,1\right)\)，\(P_2=(0,1)\)，\(P_3=\left(\frac14,\frac12\right)\). 又
\(y_4=\frac{y_1+y_2}{2}=1\)，\(y_5=\frac{y_2+y_3}{2}=\frac34\). 因此
\(a_1=\frac12y_1+y_2+y_3=2\)，\(a_2=\frac12y_2+y_3+y_4=2\)，\(a_3=\frac12y_3+y_4+y_5=2\).
由递推关系 \(y_{n+3}=\frac{y_n+y_{n+1}}2\)，有
\[
a_{n+1}=\frac12y_{n+1}+y_{n+2}+y_{n+3}
=\frac12y_n+y_{n+1}+y_{n+2}=a_n
\]
所以 \(a_n\) 是常数列，且 \(a_n=2\).
\item 由 \(a_n=2\) 得
\(y_{n+1}+y_{n+2}=2-\frac12y_n\). 另一方面，
\(y_{n+4}=\frac{y_{n+1}+y_{n+2}}2\)，所以
\(y_{n+4}=1-\frac14y_n\)，\(n\in\N^*\).
\item 由上式，
\[
b_{n+1}=y_{4n+8}-y_{4n+4}
=\left(1-\frac14y_{4n+4}\right)-\left(1-\frac14y_{4n}\right)
=-\frac14\left(y_{4n+4}-y_{4n}\right)
=-\frac14b_n
\]
又 \(y_4=1\)，\(y_5=y_6=\frac34\)，\(y_7=\frac78\)，\(y_8=\frac34\)，故
\(b_1=y_8-y_4=-\frac14\ne0\). 因此 \(\{b_n\}\) 是等比数列，公比为 \(-\frac14\).
\end{enumerate}
\end{solution}
